\subsection{某些线性群的李代数}

设 \(\mathbf{E}\) 为完备可赋范空间。于是 \(\mathcal{L}(\mathbf{E})\) 是含单位的完备可赋范代数，而 \(\mathbf{GL}(\mathbf{E})\) 是李群。由第 9 节命题 33 的推论，若把 \(\mathrm{T}_{1}(\mathbf{GL}(\mathbf{E}))\) 典范地与 \(\mathcal{L}(\mathbf{E})\) 认同，则 \(\mathrm{L}(\mathbf{GL}(\mathbf{E}))\) 上的李代数结构由括号 \((x,y)\mapsto xy - yx\)（\(\mathcal{L}(\mathbf{E})\) 的两个元素的括号）给出。特别地，\(\mathrm{L}(\mathbf{GL}(n,\mathbf{K}))\) 典范地与 \(\mathrm{gl}(n,\mathbf{K})\) 认同 (Chapter I, §1, no. 2)。

命题 35. 设 E 为有限维向量空间。设 \(\phi\) 为态射 \(g \mapsto \det g\)（从李群 \(\mathbf{GL}(E)\) 到李群 \(\mathbf{K}^*\) 中）。\(\mathrm{L}(\phi)\)（从 \(\mathcal{L}(E)\) 到 \(\mathbf{K}\)）是映射 \(x \mapsto \mathrm{Tr}x\)。\(\mathbf{SL}(E)\)（\(\phi\) 的核）是 \(\mathbf{GL}(E)\) 的李子群，其李代数为 \(\mathfrak{sl}(E)\)。

我们选取 E 的一个范数和一个基。行列式的展开式证明：

\[
\det (1 + u) \in 1 + \operatorname{Tr} u + o (\| u \|)
\]

当 \(u\) 在 \(\mathcal{L}(\mathrm{E})\) 中趋于 0 时成立。因而，对 \(x \in \mathcal{L}(\mathrm{E}) = \mathrm{L}(\mathbf{GL}(\mathrm{E}))\)，利用第 9 节的命题 34，得

\[
\mathrm{L} (\phi) (x) = \langle x, \phi \rangle = \operatorname{Tr} x.
\]

由此可得 \(\phi\) 是浸没。因此，\(\operatorname{Ker} \phi = \mathbf{SL}(E)\) 是 \(\mathbf{GL}(E)\) 的李子群，其李代数为 \(\operatorname{Ker} L(\phi) = \mathfrak{sl}(E)\)。

设 \(\mathbf{E}_1, \ldots, \mathbf{E}_n\) 为完备可赋范空间，\(\mathbf{E}\) 为它们的直和。每个 \(x \in \mathcal{L}(\mathbf{E})\) 都可以表示为一个矩阵 \((x_{ij})_{1 \leqslant i, j \leqslant n}\)，其中 \(x_{ij} \in \mathcal{L}(\mathbf{E}_i, \mathbf{E}_j)\)。

命题 36. 设 \(\mathbf{I}\) 为 \(\{1,2,\ldots,n\}\) 的子集，\(\mathbf{G}\) 为 \(\mathbf{GL}(\mathbf{E})\) 中由 \(g = (g_{ij})_{1\leqslant i,j\leqslant n}\in \mathbf{GL}(\mathbf{E})\)（满足 \(g_{ij} = 0\)（当 \(i < j\)）且 \(g_{ii} = 1\)（当 \(i\in \mathbf{I}\)））组成的子群。则 \(\mathbf{G}\) 是 \(\mathbf{GL}(\mathbf{E})\) 的李子群，且 \(\mathbf{L}(\mathbf{G})\) 是由 \(x = (x_{ij})_{1\leqslant i,j\leqslant n}\in \mathcal{L}(\mathbf{E})\)（满足 \(x_{ij} = 0\)（当 \(i < j\)）且 \(x_{ii} = 0\)（当 \(i\in \mathbf{I}\)））组成的集合。设 \(\mathbf{S}\) 为由 \((x_{ij})\in \mathcal{L}(\mathbf{E})\)（满足 \(x_{ij} = 0\)（当 \(i < j\)）且 \(x_{ii} = 0\)（当 \(i\in \mathbf{I}\)））组成的集合。则 \(\mathbf{G}\) 是 \(\mathbf{GL}(\mathbf{E})\) 与仿射子空间 \(1 + \mathbf{S}\)（\(\mathcal{L}(\mathbf{E})\) 的）的交。因而 \(\mathbf{G}\) 是 \(\mathbf{GL}(\mathbf{E})\) 的子流形，且 \(\mathbf{G}\) 在 1 处的切空间与 \(\mathbf{S}\) 认同。

特别地，在 \(\mathbf{GL}(n,\mathbf{K})\) 中，按 Integration, Chapter VII, § 3, no. 3 定义的 全下三角形子群与严格下三角形子群是分别以 \(t(n,\mathbf{K})\) 与 \(n(n,\mathbf{K})\) 为李代数的李子群 (Chapter I, § 1, no. 2)。

命题 37. 设 A 为完备可赋范含单位结合代数，\(x \mapsto x^t\) 为 A 到 A 中的连续线性映射，满足 \((x^t)^t = x\) 且 \((xy)^t = y^t x^t\)（对所有 \(x, y\)，A 中的元素）。设 K 的特征 \(\neq 2\)。设 G 为 A* 中由 \(x \in A\)（满足 \(xx^t = x^t x = 1\)）组成的子群。则 G 是 A* 的李子群，且 L(G) 是由 \(y \in A\)（满足 \(y^t = -y\)）组成的集合。

设 \(S\)（相应地 \(S'\)）为 \(y \in A\) 中满足 \(y = y^t\)（相应地 \(y = -y^t\)）者组成的集合。则 \(S, S'\) 是 \(A\) 的闭向量子空间。公式

\[
y = \frac {1}{2} (y + y ^ {t}) + \frac {1}{2} (y - y ^ {t})
\]

证明 A 是 S 与 \(S'\) 的拓扑直和。设 f 为 A 到 S 中由 \(f(x) = xx^{t}\) 定义的映射。这个映射是解析的。对所有 \(y \in A\)，\(f(1 + y) = 1 + y + y^{t} + yy^{t}\)；在 A 上选取一个与其代数结构相容的范数；则

\[
f (1 + y) \in 1 + y + y ^ {t} + o (\| y \|) \quad \text { for   } y \text {   tending   to   } 0.
\]

于是，\(\mathrm{T}_{1}(f)(y)=y+y^{t}\)，从而 f 在 1 处是浸没。因此，存在 A 中 1 的开邻域 U，使得 \(U\cap G\) 是 U 的子流形。因而 (§ 1, no. 3, Proposition 6) G 是 \(A^{*}\) 的李子群。此外，\(\mathrm{L}(G)=\mathrm{T}_{e}(G)=\mathrm{Ker}\;\mathrm{T}_{1}(f)\)。

推论 1. 设 K 的特征 \(\neq 2\)。设 E 为 K 上的有限维向量空间，\(\phi\) 为 E 上的非退化对称（相应地交错）双线性型。对所有 \(u \in \mathcal{L}(\mathrm{E})\)，用 \(u^*\) 表示 \(u\) 关于 \(\phi\) 的伴随。设 G 为 \(\phi\) 的正交群（相应地辛群）。则 G 是 \(\mathbf{GL}(\mathrm{E})\) 的李子群，且 \(\mathrm{L}(\mathrm{G})\) 是由 \(x \in \mathcal{L}(\mathrm{E})\)（满足 \(x^* = -x\)）组成的集合。

我们对 A = L(E) 且 \(x^{t} = x^{*}\) 的情形应用命题 37。

注。设 B 为 E 的一个基，J 为 \(\phi\) 关于 B 的矩阵。则 \(\mathrm{L}(G)\) 是 \(\mathcal{L}(\mathrm{E})\) 中其关于 B 的矩阵 X 满足方程

\[
{ } ^ { t } X = - J X J ^ { - 1 } .
\]

这由 Algebra, Chapter IX, § 1, formula (50) 得出。

推论 2. 设 E 为复（相应地实）希尔伯特空间，U 为 E 的酉群。则 U 是 \(\mathbf{GL}(\mathrm{E})\) 的实李子群，且 \(\mathbf{L}(\mathbf{U})\) 是由 \(x \in \mathcal{L}(\mathrm{E})\)（满足 \(x^{*} = -x\)）组成的集合。

我们把 \(\mathbf{A} = \mathcal{L}(\mathbf{E})\) 看作 \(\mathbf{R}\) 上的代数并取 \(x^t = x^*\)，应用命题 37。

推论 3. 设 E 为有限维复向量空间，\(\phi\) 为 E 上的非退化厄米特半双线性型，U 为 \(\phi\) 的酉群。则 U 是 \(\mathbf{GL}(\mathbf{E})\) 的实李子群，且 \(\mathbf{L}(\mathbf{U})\) 是由使 ix 为厄米特的那些 \(x \in \mathcal{L}(\mathbf{E})\) 组成的集合。

当 \(E \neq \{0\}\) 时，U 不是复李群 \(\mathbf{GL}(E)\) 的李子群，因为 \(\mathbf{L}(U)\) 不是 \(\mathcal{L}(E)\) 的复向量子空间。

